\subsection{State Space of a Single-Qubit System}

The \definition{State Space} of a physical system is the set of all possible states that the system can occupy. For a \textbf{single qubit}, the state space consists of all vectors of the form:
\begin{equation*}
    \ket{\psi} = a\ket{0} + b\ket{1}
\end{equation*}
\begin{itemize}
    \item $a, b \in \mathbb{C}$ are \textbf{complex probability amplitudes}.
    \item $\ket{0}$ and $\ket{1}$ form an \textbf{orthonormal basis} (the computational basis).
    \item The \textbf{normalization condition} must hold: $\left|a\right|^2 + \left|b\right|^2 = 1$. This ensures that the total probability of all measurement outcomes is equal to 1.
\end{itemize}
\begin{flushleft}
    \textcolor{Green3}{\faIcon{book} \textbf{Degrees of Freedom of a Qubit}}
\end{flushleft}
At first glance, the coefficients $a$ and $b$ appear to provide \textbf{4 real parameters}:
\begin{gather*}
    a = a_r + i a_i \quad \Rightarrow \quad a_r, a_i \quad \text{(2 parameters)} \\
    b = b_r + i b_i \quad \Rightarrow \quad b_r, b_i \quad \text{(2 parameters)}
\end{gather*}
So initially, it seems that a qubit has \textbf{four degrees of freedom}. However, not all of these correspond to physically distinguishable states.
\begin{itemize}
    \item \important{Normalization constraint}. Quantum states must be normalized, which imposes the condition:
    \begin{equation*}
        \left|a\right|^2 + \left|b\right|^2 = a_r^2 + a_i^2 + b_r^2 + b_i^2 = 1
    \end{equation*}
    Geometrically, it defines the surface of a \textbf{unit sphere in four dimension} (also known as a 3-sphere, $S^3$). This constraint effectively removes \textbf{one degree of freedom}, since the amplitudes must lie on the surface of a 3-sphere in 4-dimensional space.

    \begin{figure}[!htp]
        \centering
        \resizebox{\linewidth}{!}{
            \begin{tikzpicture}[>=Latex]
                % --- S1 ---
                \begin{scope}[shift={(-5,0)}]
                    \draw[->] (-1.4,0) -- (1.4,0) node[below] {$x_1$};
                    \draw[->] (0,-1.4) -- (0,1.4) node[left] {$x_2$};
                    \draw[thick, blue] (0,0) circle (1);
                    \node at (0,-2.0) {$S^1 \subset \mathbb{R}^2$};
                    \node at (0,-2.6) {\small $x_1^2+x_2^2=1$};
                \end{scope}

                % --- S2 ---
                \begin{scope}[shift={(0,0)}]
                    \draw[->] (-1.5,0) -- (1.5,0) node[below] {$x_1$};
                    \draw[->] (0,-1.5) -- (0,1.5) node[left] {$x_3$};
                    \draw[dashed] (0,0) ellipse (1 and 0.35);
                    \draw[thick, blue] (0,0) circle (1);
                    \node at (0,-2.0) {$S^2 \subset \mathbb{R}^3$};
                    \node at (0,-2.6) {\small $x_1^2+x_2^2+x_3^2=1$};
                \end{scope}

                % --- S3 ---
                \begin{scope}[shift={(5,0)}]
                    \draw[thick, blue] (0,0) circle (1);
                    \draw[dashed] (0,0) ellipse (1 and 0.35);
                    \node at (0,1.8) {\small not directly drawable};
                    \node at (0,-2.0) {$S^3 \subset \mathbb{R}^4$};
                    \node at (0,-2.6) {\small $x_1^2+x_2^2+x_3^2+x_4^2=1$};
                \end{scope}

                % arrows
                \draw[->, thick] (-3.2,0) -- (-1.8,0);
                \draw[->, thick] (1.8,0) -- (3.2,0);
            \end{tikzpicture}
        }
        \caption{Analogy between the unit circle $S^1$, the ordinary sphere $S^2$, and the 3-sphere $S^3$. The normalized qubit state lives on $S^3$ because it is specified by four real coordinates constrained by one equation. In this figure, we can only visualize $S^1$ and $S^2$, while $S^3$ cannot be directly drawn in three-dimensional space, so we represent it with a dashed ellipse to indicate that it is a higher-dimensional object.}
        \label{fig:s1-s2-s3-analogy}
    \end{figure}


    \newpage
    
    
    \item \important{Global Phase Equivalence}. Two states that differ only by a \textbf{global phase} represent the \textbf{same physical qubit}:
    \begin{equation*}
        \ket{\psi'} = e^{i\phi}\ket{\psi}
    \end{equation*}
    For example:
    \begin{equation*}
        \ket{\psi} = \frac{1}{\sqrt{2}}\left(\ket{0} + \ket{1}\right)
        \quad
        e^{i \pi / 3} \ket{\psi} = \frac{e^{i \pi / 3}}{\sqrt{2}}\left(\ket{0} + \ket{1}\right)
    \end{equation*}
    \textcolor{Green3}{\faIcon{question-circle} \textbf{Why is it physically irrelevant?}} Measurement probabilities are given by:
    \begin{equation*}
        P(i) = \left|\braket{i}{\psi}\right|^2
    \end{equation*}
    If we multiply the state by a global phase:
    \begin{equation*}
        \left|\braket{i}{e^{i\gamma}\psi}\right|^{2} =
        \left|e^{i\gamma} \braket{i}{\psi}\right|^{2} =
        \left|\braket{i}{\psi}\right|^{2}
    \end{equation*}
    Because $\left|e^{i\gamma}\right| = 1$ for any real $\gamma$.

    \begin{proof}[Mathematical proof that $\left|e^{i\gamma}\right| = 1$]
        We can express $e^{i\gamma}$ using Euler's formula:
        \begin{equation*}
            e^{i\gamma} = \cos(\gamma) + i\sin(\gamma)
        \end{equation*}
        The magnitude of a complex number $z = x + iy$ is given by:
        \begin{equation*}
            \left|z\right| = \sqrt{x^2 + y^2}
        \end{equation*}
        Applying this to $e^{i\gamma}$:
        \begin{align*}
            \left|e^{i\gamma}\right| &= \sqrt{\left(\cos(\gamma)\right)^2 + \left(\sin(\gamma)\right)^2} \\[.2em]
            &\downarrow \text{using the trigonometric identity} \\[.2em]
            &= \sqrt{\cos^2 \gamma + \sin^2 \gamma} \\[.2em]
            &= \sqrt{1} \\[.2em]
            &= 1
        \end{align*}
    \end{proof}
    Therefore, \textbf{no experiment can distinguish between} $\ket{\psi}$ \textbf{and} $e^{i\phi}\ket{\psi}$, since \textbf{they yield the same measurement probabilities} (due to $\left|e^{i\phi}\right| = 1$).

    \textcolor{Green3}{\faIcon{question-circle} \textbf{How does this remove a degree of freedom?}} Unlike normalization, the global phase does \textbf{not} impose a constrain equation. Instead, it introduces a \textbf{redundancy} in our description:
    \begin{itemize}
        \item Many different vectors in $S^{3}$ correspond to the \textbf{same physical state}.
        \item All vectors of the norm $e^{i\gamma}\ket{\psi}$ are considered \textbf{physically equivalent} to $\ket{\psi}$.
    \end{itemize}
    Mathematically, we are \textbf{identifying} all these vectors as a single point. This process is called forming a \textbf{quotient space}: the set of physically distinct pure qubit states is $S^{3} \setminus S^{1} \cong S^{2}$, which is a 2-dimensional surface (the Bloch sphere, see \autoref{fig:s1-s2-s3-analogy}, \autopageref{fig:s1-s2-s3-analogy}).
\end{itemize}
After accounting for these constraints, we find that a single qubit state is fully described by \textbf{two real parameters} (two degrees of freedom). Therefore, any qubit state can be written in the form:
\begin{enumerate}
    \item A normalized qubit can be written as:
    \begin{equation*}
        \ket{\psi} = a\ket{0} + b\ket{1} \qquad a,b \in \mathbb{C} \quad \left|a\right|^{2} + \left|b\right|^{2} = 1
    \end{equation*}
    
    
    \item Any complex number can be written using its \textbf{magnitude} and \textbf{phase}:
    \begin{equation*}
        a = \left|a\right| e^{i\alpha} \qquad b = \left|b\right| e^{i\beta}
    \end{equation*}
    Where $\left|a\right|$ and $\left|b\right|$ are non-negative real numbers, and $\alpha$ and $\beta$ are real numbers representing the phases. Substituting these into the expression for $\ket{\psi}$ gives:
    \begin{equation*}
        \ket{\psi} = \underbrace{\left|a\right| e^{i\alpha}}_{a} \ket{0} + \underbrace{\left|b\right| e^{i\beta}}_{b} \ket{1}
    \end{equation*}
    
    
    \item We can factor out the common global phase $e^{i\alpha}$:
    \begin{equation*}
        \ket{\psi} = e^{i\alpha} \left( \left|a\right| \ket{0} + \left|b\right| e^{i(\beta - \alpha)} \ket{1} \right)
    \end{equation*}
    Since the global phase $e^{i\alpha}$ has \textbf{no physical effect}, we can omit it, leading to:
    \begin{equation*}
        \ket{\psi} = \left|a\right| \ket{0} + \left|b\right| e^{i\phi} \ket{1}
    \end{equation*}
    Where we define also the \definition{Relative Phase} $\phi = \beta - \alpha$. At this point, the state is characterized by two real magnitudes $\left|a\right|$ and $\left|b\right|$, and one relative phase $\phi$.


    \item From normalization:
    \begin{equation*}
        \left|a\right|^{2} + \left|b\right|^{2} = 1
    \end{equation*}
    We recognize a familiar trigonometric identity:
    \begin{equation*}
        \cos^2 \left(x\right) + \sin^2 \left(x\right) = 1
    \end{equation*}
    Therefore, we can parametrize the magnitudes using a single angle $\theta$:
    \begin{equation*}
        \left|a\right| = \cos\left(\frac{\theta}{2}\right) \qquad \left|b\right| = \sin\left(\frac{\theta}{2}\right)
    \end{equation*}
    With $0 \leq \theta \leq \pi$ to ensure that $\left|a\right|$ and $\left|b\right|$ are non-negative. The reason for using $\frac{\theta}{2}$ instead of $\theta$ will become clear when we discuss the Bloch sphere representation (\autoref{fig:bloch-sphere}, \autopageref{fig:bloch-sphere}).


    \item Substituting the parametrized magnitudes into the state:
    \begin{equation}\label{eq:canonical-parametrization}
        \ket{\psi} = \underbrace{\cos\left(\frac{\theta}{2}\right)}_{\left|a\right|}\ket{0} + e^{i\phi}\underbrace{\sin\left(\frac{\theta}{2}\right)}_{\left|b\right|} \ket{1}
    \end{equation}
    This is the \textbf{canonical parametrization} of a single qubit state.
\end{enumerate}
Regardless of the specific values of $\theta$ and $\phi$ in \autoref{eq:canonical-parametrization}:
\begin{itemize}
    \item $\theta$: controls the relative weights of $\ket{0}$ and $\ket{1}$ in the superposition, determining the probabilities of measuring each state.
    \item $\phi$: controls the relative phase between the $\ket{0}$ and $\ket{1}$ components. While it does not affect measurement probabilities, it is crucial for interference effects and the behavior of the qubit under quantum gates (which we will discuss later).
\end{itemize}

\begin{flushleft}
    \textcolor{Green3}{\faIcon{book} \textbf{Bloch Sphere Representation of a Qubit}}\index{Bloch Sphere Representation of a Qubit}
\end{flushleft}
After normalization and removal of the physically irrelevant global phase, a pure single-qubit state depends on only two real parameters. Therefore, it is natural to represent it geometrically as a point on a two-dimensional surface. This surface is the \textbf{Bloch sphere}.

\highspace
Any pure qubit can be written in the form:
\begin{equation*}
    \ket{\psi}
    =
    \cos\left(\frac{\theta}{2}\right)\ket{0}
    +
    e^{i\phi}\sin\left(\frac{\theta}{2}\right)\ket{1},
\end{equation*}
Where:
\begin{itemize}
    \item $\theta \in \left[0,\pi\right]$ is the \textbf{polar angle},
    \item $\phi \in \left[0,2\pi\right)$ is the \definition{Azimuthal Angle}, representing the \textbf{relative phase}.
\end{itemize}
This expression gives a one-to-one correspondence between pure qubit states and points on the surface of a unit sphere:
\begin{itemize}
    \item $\ket{0}$ corresponds to the north pole, since $\theta = 0$:
    \begin{equation*}
        \ket{0} = \cos\left(\dfrac{0}{2}\right)\ket{0} + e^{i\phi}\sin\left(\dfrac{0}{2}\right)\ket{1} = 1 \cdot \ket{0} + \underbrace{e^{-i}}_{\phi = 0 - 1} \cdot 0 \cdot \ket{1} = \ket{0}
    \end{equation*}
    \item $\ket{1}$ corresponds to the south pole, since $\theta = \pi$:
    \begin{equation*}
        \ket{1} = \cos\left(\dfrac{\pi}{2}\right)\ket{0} + e^{i\phi}\sin\left(\dfrac{\pi}{2}\right)\ket{1} = 0 \cdot \ket{0} + \underbrace{e^{i}}_{\phi = 1 - 0} \cdot 1 \cdot \ket{1} = e^{i} \ket{1} \Rightarrow \ket{1}
    \end{equation*}
    \textcolor{Red2}{\faIcon{exclamation-triangle} \textbf{Attention!}} $e^{i}\ket{1}$ becomes $\ket{1}$ after removing the global phase, so $\ket{1}$ is indeed at the south pole.
    \item States on the equator correspond to equal-weight superpositions of $\ket{0}$ and $\ket{1}$.
\end{itemize}
In particular:
\begin{align*}
    \text{Hadamard} & \quad \ket{+}  = \frac{1}{\sqrt{2}}(\ket{0}+\ket{1}), \quad \ket{-} = \frac{1}{\sqrt{2}}(\ket{0}-\ket{1}),\\[.5em]
    \text{Imaginary} & \quad \ket{i} = \frac{1}{\sqrt{2}}(\ket{0}+i\ket{1}), \quad \ket{-i} = \frac{1}{\sqrt{2}}(\ket{0}-i\ket{1}),
\end{align*}
All lie on the equator and differ only by their relative phase.

\highspace
The Bloch sphere provides an intuitive visualization of a qubit: moving along the vertical direction changes the relative weights of $\ket{0}$ and $\ket{1}$, while rotating around the vertical axis changes the relative phase between them.
\begin{figure}[!htp]
    \centering
    \tdplotsetmaincoords{72}{120}
    \begin{tikzpicture}[tdplot_main_coords, scale=3.3, >=Latex]

        \def\r{1}

        % Axes: positive directions
        \draw[->, thick] (0,0,0) -- (1.25,0.05,0) node[anchor=north east] {$x$};
        \draw[->, thick] (0,0,0) -- (0,1.25,0) node[anchor=west] {$y$};
        \draw[->, thick] (0,0,0) -- (0,0,1.25) node[anchor=south] {$z$};

        % Axes: negative directions (lighter dashed hints)
        \draw[dashed, thin, black!55] (0,0,0) -- (-1.15,0,0);
        \draw[dashed, thin, black!55] (0,0,0) -- (0,-1.15,0);
        \draw[dashed, thin, black!55] (0,0,0) -- (0,0,-1.15);

        % Sphere
        \tdplotdrawarc[dashed]{(0,0,0)}{\r}{0}{180}{}{}
        \tdplotdrawarc[thick]{(0,0,0)}{\r}{180}{360}{}{}

        \tdplotsetrotatedcoords{90}{0}{90}
        \tdplotdrawarc[tdplot_rotated_coords,dashed]{(0,0,0)}{\r}{0}{180}{}{}
        \tdplotdrawarc[tdplot_rotated_coords,thick]{(0,0,0)}{\r}{180}{360}{}{}

        \tdplotsetrotatedcoords{0}{90}{90}
        \tdplotdrawarc[tdplot_rotated_coords,dashed]{(0,0,0)}{\r}{0}{180}{}{}
        \tdplotdrawarc[tdplot_rotated_coords,thick]{(0,0,0)}{\r}{180}{360}{}{}
        \tdplotsetrotatedcoords{0}{0}{0}

        % States
        \fill (0,0,1) circle (0.015);
        \node[anchor=south east] at (0,0,1.08) {$\ket{0}$};

        \fill (0,0,-1) circle (0.015);
        \node[anchor=north east] at (0,0,-1.08) {$\ket{1}$};

        \fill (1,0,0) circle (0.015);
        \node[anchor=north] at (1.11,0.10,0) {$\ket{+}$};

        \fill (-1,0,0) circle (0.015);
        \node[anchor=south] at (-1.11,0,0) {$\ket{-}$};

        \fill (0,1,0) circle (0.015);
        \node[anchor=south west] at (0.05,1.12,0) {$\ket{i}$};

        \fill (0,-1,0) circle (0.015);
        \node[anchor=north east] at (0.01,-1.10,0.05) {$\ket{-i}$};

        % Generic vector
        \def\thetavec{52}
        \def\phivec{38}
        \tdplotsetcoord{P}{\r}{\thetavec}{\phivec}

        \pgfmathsetmacro{\px}{sin(\thetavec)*cos(\phivec)}
        \pgfmathsetmacro{\py}{sin(\thetavec)*sin(\phivec)}

        % Bloch vector
        \draw[->, very thick, red] (0,0,0) -- (P);
        \node[anchor=south west, red] at (0, 0.12, 0.16) {$\ket{v}$};

        % Projection helpers:
        % vertical dashed red line
        \draw[dashed, red!70] (P) -- (\px,\py,0);
        % red line in equatorial plane from origin to the foot of the vertical
        \draw[dashed, red!70] (0,0,0) -- (\px,\py,0);

        % Angles
        \tdplotsetthetaplanecoords{\phivec}
        \tdplotdrawarc[tdplot_rotated_coords,->]
            {(0,0,0)}{0.50}{0}{\thetavec}
            {anchor=south west}{$\theta$}

        \tdplotsetrotatedcoords{0}{0}{0}
        \tdplotdrawarc[->]
            {(0,0,0)}{0.20}{0}{\phivec}
            {anchor=north east}{$\phi$}

    \end{tikzpicture}
    \caption{Bloch sphere representation of a pure qubit state.}
    \label{fig:bloch-sphere}
\end{figure}